This thesis established an expression-guided, multi-evidence framework for
evaluating seven cartilage- and growth-plate-relevant SLRPs across
representative vertebrates. BGN, DCN, FMOD, PRELP, EPYC, LUM, and OGN generally
retain a strongly conserved secreted LRR-rich protein core, and most assignments
are supported jointly by sequence alignment, domain architecture, gene-specific
phylogenetic placement, representative coding-exon organization, and genomic
neighbourhood. The seven completed SynVoy reports yielded a 98-row evidence
table, although 25 priority loci and subsequent rapid confirmations still
require manual biological review. BGN and FMOD are the strongest combined
biological and technical anchors; DCN is both a biologically supported candidate
and the essential class-I paralog control; PRELP, EPYC, and LUM remain well-
supported principal candidates; and OGN is most appropriately retained as an
exploratory vertebrate member.

Conservation was not uniform and did not make the comparative analysis
uninformative. The workflow exposed an ASPN-like chicken BGN-locus product,
excluded an amphioxus SLRP-like sequence from the confident OGN set, and retained
spotted-gar OGN, whale-shark EPYC, and lamprey LUM with explicit caveats. Stable
coding-exon counts coexisted with substantial intron-driven variation in locus
span, and mouse and rat growth-plate expression were not identical. These cases
show why orthology should be judged from concordant evidence rather than from a
gene label, one BLAST hit, or one conserved domain.

The evolutionary contribution is therefore not a claim that one sampled living
species represents the origin of an SLRP gene. Literature places an SLRP-like
precursor early in chordates and family expansion in early vertebrates. The
amphioxus--lamprey--cartilaginous-fish--bony-vertebrate bracket tested where
modern gene identities become confidently recognizable. The results support a
broadly established jawed-vertebrate SLRP repertoire while revealing decreasing
assignment certainty in the deepest lineages and in incomplete shark models.

Supplementary protein motifs supported every retained gene assignment, whereas
comparative promoter analysis did not reveal a corrected shared targeted
regulatory signature. Curated phenotype databases strengthened the skeletal
context of BGN, FMOD, and EPYC and provided broader ECM context for the other
genes, but did not justify ranking genes by annotation counts. Overall, the
results support evolutionary constraint on an extracellular-matrix SLRP toolkit
relevant to juvenile cartilage biology while clearly separating conserved
molecular compatibility from experimentally demonstrated growth-plate function
in each species. Targeted manual review of the prioritized ambiguous loci will
finalize the locus-level evidence layer.
